\article{正规化因子、缓增表示与 $L$ 群}{Normalizing factors, tempered representations, and $L$-groups}{A. W. Knapp 与 Gregg Zuckerman}{1}{93--105}
几位讲演者提到了与主系列表示（principal series representations）相关的交织算子（intertwining operators），特别是 $\SL(2,\R)$ 的情形。这些算子及其正规化（normalizations）在迹公式（trace formula）、主系列可约性问题，以及使某些表示成为酉表示的内积构造中都有重要作用。我们将在一个线性实约化群（linear real reductive group）$G$ 的框架中回顾这些算子及其正规化的发展；$G$ 的定义同 Wallach 的讲座，并假设其单位连通分支具有紧中心。\footnote{两位作者均获美国国家科学基金会资助。原文 AMS（MOS）学科分类（1970）：主要 22E45。}

我们尤其关心它们在酉主系列以及来自抛物子群的其他连续系列表示（continuous series of representations）的可约性问题中的应用，因为这些问题的答案导出不可约缓增表示（irreducible tempered representations）的分类，从而补充 Langlands 分类 [12]。第2节借助三个易于计算的有限群 $W,W',R$，介绍相应的可约性判据。

作者在1975--76年的讲座中曾指出，在某些特殊情形下，最初用根定义的 $W,W',R$ 可以用余根（co-roots）定义。在此基础上，Langlands [14] 得以在一般情形下用 $L$ 群（$L$-group）重新定义这些群。第3节给出他的定义与一个例子。第4节概述利用 $R$ 群分类不可约缓增表示的早期工作 [8]、[9]。

\origsection{1}{交织算子与正规化因子（Intertwining operators and normalizing factors）}
在 $G$ 中固定极大紧子群（maximal compact subgroup）$K$ 及 Cartan 对合（Cartan involution）$\theta$。对每个抛物子群 $P$，取 Langlands 分解 $P=MAN$，其中 $MA$ 在 $\theta$ 下稳定，$M$ 为单位连通分支具有紧中心的线性实约化群。设 $\xi$ 为 $M$ 的不可约酉表示，$\nu$ 为 $A$ 的 Lie 代数 $\mathfrak a$ 上的复值线性泛函；对这对数据定义表示 $U_P(\xi,\nu,x)$（$x\in G$）为
\begin{equation}
 U_P(\xi,\nu)=\operatorname{ind}_{MAN}^G(\xi\otimes e^\nu\otimes1).
 \tag{1.1}
\end{equation}
我们约定 $G$ 在诱导表示中从左作用。表示空间的元素 $f$ 满足
\[
 f(xman)=\exp\bigl(- (\nu+\rho_P)\log a\bigr)\xi(m)^{-1}f(x),
\]
其中 $\rho_P$ 是与 $P$ 对应的限制根（restricted roots）的通常半和。若 $P=P_0$ 为极小抛物子群，这些表示组成非酉主系列。式 (1.1) 假定 $\xi$ 为 $M$ 的不可约酉表示，但只要不处理涉及算子伴随的公式，也允许 $\xi$ 为 $M$ 的非酉主系列表示。

对于 $\SL(2,\R)$，把非酉主系列表示空间中的函数限制到下三角群 $\overline N=\theta N=\{\bigl(\begin{smallmatrix}1&0\\ *&1\end{smallmatrix}\bigr)\}$，可在 $\R$ 上的函数空间中实现这些表示。当 $g=\bigl(\begin{smallmatrix}a&b\\c&d\end{smallmatrix}\bigr)$ 时，它们成为
\[
 \mathcal P^{\pm,\zeta}(g)f(x)=
 \begin{cases}
 |-bx+d|^{-1-\zeta}f((ax-c)/(-bx+d)),&\text{取 }+,\\
 \operatorname{sgn}(-bx+d)|-bx+d|^{-1-\zeta}f((ax-c)/(-bx+d)),&\text{取 }-.
 \end{cases}
\]
Kunze 和 Stein [10] 证明，算子
\begin{equation}
 f\longmapsto\int_{-\infty}^\infty\frac{f(x-y)\,dy}{|y|^{1-\zeta}}
 \quad\text{（对于 }\mathcal P^{+,\zeta}\text{）}
 \tag{1.2a}
\end{equation}
以及
\begin{equation}
 f\longmapsto\int_{-\infty}^\infty\frac{(\operatorname{sgn}y)f(x-y)\,dy}{|y|^{1-\zeta}}
 \quad\text{（对于 }\mathcal P^{-,\zeta}\text{）}
 \tag{1.2b}
\end{equation}
在积分收敛时，即 $\operatorname{Re}\zeta>0$ 时，交织 $\mathcal P^\zeta$ 与 $\mathcal P^{-\zeta}$。后来 [11] 他们在诱导模型中找到公式
\[
 f\longmapsto\int_{\overline N}f(gw\bar n)\,d\bar n,
 \qquad w=\begin{pmatrix}0&1\\-1&0\end{pmatrix}.
\]
这是 $f\mapsto\int_{\overline N}f(g\bar n)\,d\bar n$ 与一个相对平凡的按 $w$ 平移算子的复合。第一个算子交织从 $P=MAN$ 诱导的表示与从 $\overline P=MA\overline N$ 诱导的表示，后者在 $MA$ 上使用同样的数据。

一般地，设 $P_1=MAN_1,P_2=MAN_2$ 是具有相同 $MA$ 的两个抛物子群，定义
\begin{equation}
 A(P_2:P_1:\xi:\nu)f(x)=\int_{\overline N_1\cap N_2}f(x\bar n)\,d\bar n.
 \tag{1.3}
\end{equation}
形式上有
\begin{equation}
 U_{P_2}(\xi,\nu,g)A(P_2:P_1:\xi:\nu)
 =A(P_2:P_1:\xi:\nu)U_{P_1}(\xi,\nu,g).
 \tag{1.4}
\end{equation}
积分 (1.3) 只在一定的参数范围内收敛，其余参数处须通过解析延拓（analytic continuation）定义。为此，当 $\nu$ 变化时，需要把这些表示实现在同一个空间上。把函数限制到 $K$ 即得此空间。设 $\xi$ 作用在 $H^\xi$ 上，考虑满足 $f(km)=\xi(m)^{-1}f(k)$（$k\in K,m\in M\cap K$）的函数 $f:K\to H^\xi$。在 $g$ 作用下，将 $f(k)$ 换为 $f(g^{-1}k)$，再限制回 $K$。于是定义
\begin{equation}
 \begin{split}
 U_P(\xi,\nu,g)f(k)
 ={}&\exp\bigl(- (\rho_P+\nu)H(g^{-1}k)\bigr)\\
 &\cdot\xi(\mu(g^{-1}k))^{-1}f(\kappa(g^{-1}k)),
 \end{split}
 \tag{1.5}
\end{equation}
其中 $x$ 在 $G=KMAN$ 中分解为 $\kappa(x)\mu(x)e^{H(x)}n$。

(1.3) 的解析延拓分三步完成。第一步的关键，在对光滑紧支撑函数 $f$ 延拓 (1.2a) 时已经可以看出。由于 $f$ 紧支撑，可在有限区间上积分。将 $f(x-y)$ 在 $y=0$ 附近展开为带余项的有限 Taylor 级数。各多项式项可以直接积分，并亚纯延拓（meromorphic continuation）到整个平面。若 Taylor 级数的阶数为 $n$，余项在 $\operatorname{Re}\zeta>-n-1$ 时可积。因此整个积分亚纯延拓至全平面。

Schiffmann [15] 及 [6] 证明，当 $P_1,P_2$ 为极小抛物子群且 $A$ 为一维时，这类论证给出 (1.3) 的亚纯延拓。第二步，对 $\dim A>1$ 的极小抛物子群，Schiffmann 将 $\overline N_1\cap N_2$ 分解为子群的乘积，把 (1.3) 写成累次积分；这样解析延拓本质上归结到 $\dim A=1$。

第三步即最后一步由 Wallach [17] 及 [7] 独立得到。利用 Casselman 子表示定理（subrepresentation theorem），可将 $\xi$ 作为 $(\mathfrak m,K\cap M)$ 模嵌入 $M$ 的一个非酉主系列表示。由分步诱导定理（double induction theorem），(1.1) 是某非酉主系列的子表示，而 (1.3) 是一个非酉主系列交织算子的限制。因此，对 $K$ 有限的 $f$，(1.3) 有解析延拓；把 (1.4) 中的 $g$ 换成 $K$ 或 $G$ 的 Lie 代数 $\mathfrak g$ 的元素，等式成立。而且适当正规化 Haar 测度后，对酉的 $\xi$，在各个 $K$ 等型子空间上取伴随时，有
\begin{equation}
 A(P_2:P_1:\xi:\nu)^*=A(P_1:P_2:\xi:-\bar\nu),
 \tag{1.6}
\end{equation}
以及
\[
 A(P_2:P_1:\xi:\nu)=R(w^{-1})
 A(wP_2w^{-1}:wP_1w^{-1}:w\xi:w\nu)R(w),
\]
其中 $R(w)$ 为按 $K$ 中元素 $w$ 的右平移，$w$ 代表 Weyl 群 $W(\mathfrak a)$ 的某个元素。

这些算子在许多重要参数值处会出现极点。引入标量正规化因子，一方面消去部分极点，一方面使算子在 $\nu$ 为纯虚时成为酉算子，还使它们在复合下有良好性质。正规化因子并不唯一，不同选择适用于不同目的。

构造仍分几步。要求 $\xi$ 具有实无穷小特征（real infinitesimal character）。所需的基本引理如下 [6，第544页]，其证明使用 Weierstrass 典范乘积。
\begin{originaltheorem}{引理 1.1（Lemma 1.1）}
设 $\eta(z)$ 为平面上的亚纯函数，满足：(i) 在实轴上取实值；(ii) 在虚轴上 $\eta(z)\ge0$；(iii) 对所有 $z$ 有 $\eta(z)=\eta(-z)$。则存在平面上的亚纯函数 $\gamma(z)$，使 $\eta(z)=\gamma(-z)\gamma(z)$，且 $\gamma(z)$ 在实轴上取实值。
\end{originaltheorem}
构造的第一步处理极小抛物子群 $P=MAN$ 及实秩一（real rank one）的 $G$，即 $\dim A=1$。此时 $A(P:\overline P:\xi:\nu)A(\overline P:P:\xi:\nu)$ 是 $U_P(\xi,\nu)$ 的自交织算子。Bruhat [3] 的结果指出，对非零纯虚 $\nu$，酉表示 $U_P(\xi,\nu)$ 不可约。由此容易得到
\begin{equation}
 A(P:\overline P:\xi:\nu)A(\overline P:P:\xi:\nu)
 =\eta(\overline P:P:\xi:\nu)I,
 \tag{1.7}
\end{equation}
其中以 $\nu=z\rho_P$ 为变量时，右侧系数是 $z\in\C$ 的亚纯标量函数。可验证：
\begin{enumerate}
\item[(i)] $\nu$ 为实时 $\eta(\overline P:P:\xi:\nu)$ 为实；
\item[(ii)] $\nu$ 为纯虚时 $\eta(\overline P:P:\xi:\nu)\ge0$；
\item[(iii)] 对所有 $\nu$ 有 $\eta(\overline P:P:\xi:\nu)=\eta(\overline P:P:\xi:-\nu)$；
\item[(iv)] 该函数只依赖于 $\xi$ 的等价类；若 $\varphi$ 是保持 $K,P$ 稳定的 $G$ 自同构，则 $\eta(\overline P:P:\xi^\varphi:\nu)=\eta(\overline P:P:\xi:\nu)$。
\end{enumerate}
于是可应用引理得到 $\gamma(\overline P:P:\xi:\nu)$，并定义正规化算子
\[
 \mathcal A(\overline P:P:\xi:\nu)=\gamma(\overline P:P:\xi:\nu)^{-1}A(\overline P:P:\xi:\nu).
\]
它们满足
\[
 \mathcal A(P:\overline P:\xi:\nu)\mathcal A(\overline P:P:\xi:\nu)=I,
 \qquad
 \mathcal A(\overline P:P:\xi:\nu)^*=\mathcal A(P:\overline P:\xi:-\bar\nu).
\]
两式共同蕴含 $\mathcal A$ 在 $\nu$ 为纯虚时为酉算子。

第二步处理一般 $G$ 的极小抛物子群。利用 Schiffmann 将一般交织算子分解为本质上秩一算子的方法，以各正规化因子的乘积作为给定算子的正规化因子。然后证明
\begin{equation}
 \mathcal A(P_3:P_1:\xi:\nu)
 =\mathcal A(P_3:P_2:\xi:\nu)\mathcal A(P_2:P_1:\xi:\nu).
 \tag{1.8}
\end{equation}
第三步处理 $\dim A=1$ 的一般抛物子群 $P=MAN$。再次将 $\xi$ 作为子表示嵌入 $M$ 的某个非酉主系列。把此方法与极小抛物子群的 (1.8) 结合，可得 (1.7) 对此 $P$ 也成立，且 $\eta$ 仍具前述性质。再次应用引理得到正规化因子 $\gamma$，令 $\mathcal A=\gamma^{-1}A$。

最后一步处理一般 $G$ 的一般抛物子群，方法与一般 $G$ 的极小抛物子群相同。若相容地选择各个 $\gamma$，便得到以下结果 [7，第2460页]。
\begin{originaltheorem}{定理 1.2（Theorem 1.2）}
正规化交织算子满足：
\begin{enumerate}
\item[(i)] $\mathcal A(P_3:P_1:\xi:\nu)=\mathcal A(P_3:P_2:\xi:\nu)\mathcal A(P_2:P_1:\xi:\nu)$；
\item[(ii)] 在各个 $K$ 等型子空间上有 $\mathcal A(P_2:P_1:\xi:\nu)^*=\mathcal A(P_1:P_2:\xi:-\bar\nu)$；
\item[(iii)] 当 $\nu$ 为纯虚时，$\mathcal A(P_2:P_1:\xi:\nu)$ 为酉算子；
\item[(iv)] 若 $w\in K$ 代表 $W(\mathfrak a)$ 的元素，则
\[
 \mathcal A(P_2:P_1:\xi:\nu)=R(w^{-1})
 \mathcal A(wP_2w^{-1}:wP_1w^{-1}:w\xi:w\nu)R(w).
\]
\end{enumerate}
\end{originaltheorem}
对代表 $W(\mathfrak a)$ 元素的 $w\in K$，令
\[
 \mathcal A_P(w,\xi,\nu)=R(w)\mathcal A(w^{-1}Pw:P:\xi:\nu).
\]
由 (i)、(iv) 得到上循环关系（cocycle relation）
\begin{equation}
 \mathcal A_P(w_1w_2,\xi,\nu)=
 \mathcal A_P(w_1,w_2\xi,w_2\nu)\mathcal A_P(w_2,\xi,\nu).
 \tag{1.9}
\end{equation}
由 (ii) 得
\begin{equation}
 \mathcal A_P(w,\xi,\nu)^*=\mathcal A_P(w^{-1},w\xi,-w\bar\nu),
 \tag{1.10}
\end{equation}
因而 $\mathcal A_P(w,\xi,\nu)$ 在 $\nu$ 为纯虚时为酉算子。交织关系为
\begin{equation}
 U_P(w\xi,w\nu,g)\mathcal A_P(w,\xi,\nu)
 =\mathcal A_P(w,\xi,\nu)U_P(\xi,\nu,g),
 \tag{1.11}
\end{equation}
其中 $g$ 属于 $K$ 或 $\mathfrak g$。

应用于阿代尔（adelic）情形时，需要更具体的正规化。几乎所有位处的表示都具有 $K$ 不动向量；只有当局部算子固定这个向量时，局部正规化交织算子的无穷张量积才有定义。此条件决定正规化因子。例如，当 $P_1,P_2$ 极小且 $\xi=1$ 时，须以
\[
 \gamma(P_2:P_1:\xi:\nu)=\int_{\overline N_1\cap N_2}
 \exp\bigl(- (\rho_P+\nu)H(\bar n)\bigr)\,d\bar n
\]
正规化 (1.3)，这个积分可通过 Harish-Chandra 的 $c$ 函数计算。Schiffmann 在通向 [15] 的研究早期就进一步探索了这个想法。Langlands 据此猜想 [13，第282页]，一般情况下可用两个 $L$ 函数之商得到有效的正规化：
\begin{equation}
 \gamma(P_2:P_1:\xi:\nu)=
 \frac{L(0,\widetilde\rho_{P_2\mid P_1}\circ\varphi)}
 {L(1,\widetilde\rho_{P_2\mid P_1}\circ\varphi)}.
 \tag{1.12}
\end{equation}
这里 $\varphi$ 是对应于 $\xi\otimes e^\nu$ 的从 $\R$ 的 Weil 群（Weil group）到 $MA$ 的 $L$ 群 ${}^L(MA)$ 的同态；$\widetilde\rho_{P_2\mid P_1}$ 为 ${}^L(MA)$ 在 ${}^L\overline N_1\cap{}^LN_2$ 的 Lie 代数上表示的逆反（contragredient）。Arthur [1] 证明，只要适当正规化 Haar 测度，(1.12) 就是有效的正规化。

\origsection{2}{$R$ 群的第一种表述（The $R$-group, a first formulation）}
固定 $M$ 的离散系列表示（discrete series representation）$\xi$，考虑相应的连续系列 $U_P(\xi,\nu)$，其中 $\nu$ 为纯虚。[7] 说明，确定与 $U_P(\xi,\nu)$ 交换的算子代数 $\mathcal C(\xi,\nu)$ 的问题，如何归结为涉及某些有限群的代数问题。

固定 $\nu$。令 $W$ 为 $W(\mathfrak a)$ 中同时固定 $\nu$ 与 $\xi$ 的等价类的元素 $s$ 组成的子群。若 $w\in K$ 代表 $s\in W$，可定义 $\xi(w)$，使 $\xi$ 延拓为包含 $M$ 与 $w$ 的最小群的表示；此定义在相差一个单位根标量因子的意义下唯一。于是 $\xi(w)\mathcal A_P(w,\xi,\nu)$ 与代表元 $w$ 无关，可记为
\begin{equation}
 \xi(s)\mathcal A_P(s,\xi,\nu).
 \tag{2.1}
\end{equation}
由 (1.11) 可知，酉算子 (2.1) 将 $U_P(\xi,\nu)$ 与自身交织。下面定理的实质归于 Harish-Chandra [4]。
\begin{originaltheorem}{定理 2.1（Theorem 2.1）}
$s\in W$ 时的算子 (2.1) 张成交换算子代数（commuting algebra）$\mathcal C(\xi,\nu)$。
\end{originaltheorem}
尽管有 (1.9)，仍不能推出映射
\begin{equation}
 s\in W\longmapsto\xi(s)\mathcal A_P(s,\xi,\nu)
 \tag{2.2}
\end{equation}
为到酉算子群的同态；各 $\xi(s)$ 因子必须相容地选择，而而相容选择可能存在障碍。但 (2.1) 至少定义了到射影酉群（projective unitary group）的同态。令其核为 $W'$。
\begin{originaltheorem}{定理 2.2（Theorem 2.2）}
$W'$ 是包含在 $(\mathfrak g,\mathfrak a)$ 根集中的某个根系 $\Delta'$ 的 Weyl 群。因此，若将 $R$ 定义为 $W$ 中保持 $\Delta'$ 的正根集稳定的子群，则 $W=W'R$ 为半直积（semidirect product）。映射 (2.2) 的像由线性无关的算子组成；因此 $\mathcal C(\xi,\nu)$ 的维数等于 $R$ 的阶。\footnote{\textbf{译者注（本次校订）：}前一句的线性无关性只针对(2.2)限制到 $R$ 后得到的算子。对整个 $W$ 而言，同一 $W'$ 陪集中的算子相差非零标量，因而不可能全部线性无关；取 $R$ 作为陪集代表后，恰得到 $|R|$ 个线性无关的算子。}
\end{originaltheorem}
事实上，$R$ 的元素恰为 $W$ 中使 (2.1) 的正规化因子在 $\nu$ 处全纯的元素。因此 $W',R$ 与如下等式定义的函数 $\eta(P_2:P_1:\xi:\nu)$ 密切相关：
\[
 A(P_1:P_2:\xi:\nu)A(P_2:P_1:\xi:\nu)=\eta(P_2:P_1:\xi:\nu)I.
\]
通过分析方法，又可将这些函数联系到 $G$ 的子群的 Plancherel 测度（Plancherel measures）；这些测度已由 Harish-Chandra 确定。因此可如下具体描述 $\Delta'$。

给 $\mathfrak a$ 添上 $M$ 的 Lie 代数 $\mathfrak m$ 的一个紧 Cartan 子代数，记一般的根为 $\beta$。令 $\lambda$ 为离散系列 $\xi$ 的 Harish-Chandra 参数。设 $\alpha$ 是 $\mathfrak a$ 根，且 $2\alpha$ 不是 $\mathfrak a$ 根。若 $\alpha$ 的重数为偶数且 $2\alpha$ 不是 $\mathfrak a$ 根，定义
\[
 \mu_{\xi,\alpha}(\nu)=\prod_{\beta|_{\mathfrak a}=\alpha}\langle\lambda+\nu,\beta\rangle.
\]
若 $\alpha$ 的重数为奇数，定义
\[
 \mu_{\xi,\alpha/2}(\nu)=\mu_{\xi,\alpha}(\nu)
 =\left(\prod_{\beta|_{\mathfrak a}=c\alpha,\ c>0}\langle\lambda+\nu,\beta\rangle\right)f_{\xi,\alpha}(\nu),
\]
其中 $f_{\xi,\alpha}(\nu)$ 取
\[
 \tan\bigl(\pi\langle\nu,\alpha\rangle/|\alpha|^2\bigr)
 \quad\text{或}\quad
 \cot\bigl(\pi\langle\nu,\alpha\rangle/|\alpha|^2\bigr),
\]
分别对应于
\[
 \xi(\gamma_\alpha)=-(-1)^{2\langle\rho_\alpha,\alpha\rangle/|\alpha|^2}I
 \quad\text{或}\quad
 +(-1)^{2\langle\rho_\alpha,\alpha\rangle/|\alpha|^2}I.
\]
这里 $\alpha$ 是实根（real root），存在相应的 $\SL(2,\R)\to G$ 同态，$\gamma_\alpha$ 表示 $-I$ 的像。

现在可将 $\Delta'$ 刻画为
\[
 \Delta'=\{\mathfrak a\text{ 的根 }\alpha:s_\alpha\in W\text{ 且 }\mu_{\xi,\alpha}(\nu)=0\},
\]
其中 $s_\alpha$ 表示与 $\alpha$ 对应的反射（reflection）。同前，
\[
 R=\{p\in W:p\alpha>0\text{ 对 }\Delta'\text{ 中每个 }\alpha>0\text{ 成立}\}.
\]
于是有以下结果 [5]、[8]。
\begin{originaltheorem}{定理 2.3（Theorem 2.3）}
(a) $R$ 是二元群 $\Z_2$ 的有限直和。

(b) 对 $s\in R$，可选择算子 $\xi(s)$ 使 (2.2) 为到酉算子群的同态；因此代数 $\mathcal C(\xi,\nu)$ 可交换。

(c) 存在正交的正实根集合 $\mathcal H=\{\alpha_1,\ldots,\alpha_q\}$，满足：(i) $\mathcal H$ 的线性包中的根只有 $\pm\alpha_j$；(ii) 每个 $r\in R$ 均形如 $s_{\alpha_{j_1}}\cdots s_{\alpha_{j_n}}$；(iii) 每个 $\alpha_j$ 都出现在某个 $r\in R$ 的分解中。
\end{originaltheorem}

\origsection{3}{$R$ 群的第二种表述（The $R$-group, a second formulation）}
$G$ 的一个容许表示若其 $K$ 有限矩阵系数对每个 $\varepsilon>0$ 均属于 $L^{2+\varepsilon}(G)$，便称为缓增。已知每个不可约缓增表示都是第2节所考虑类型的某个表示的直和项，即从 $M$ 上的离散系列及 $A$ 上的酉特征诱导的表示中的直和项。$R$ 群可用于分解这些诱导表示，并给出不可约缓增表示的分类；第4节将再谈这一点。

现在换一个角度考察 $W,W',R$，先从 $G$ 分裂且抛物子群极小的情形说起。此时 $M$ 为有限阿贝尔群，由对应于各根的 $\SL(2,\R)$ 中 $\bigl(\begin{smallmatrix}-1&0\\0&-1\end{smallmatrix}\bigr)$ 的像 $\gamma_\alpha$ 生成。这里 $\xi$ 是 $M$ 的特征，设 $\nu=0$。根系记为 $\Delta$。容易看出
\[
 \Delta'=\{\alpha\in\Delta:\xi(\gamma_\alpha)=1\}.
\]
$\gamma_\alpha$ 在 Weyl 群元素下的变换与可能的预期不同。正确规则为
\[
 s_\beta\gamma_\alpha s_\beta^{-1}=\gamma_{s_\beta\alpha}
 =\gamma_\alpha\gamma_\beta^{2\langle\beta,\alpha\rangle/|\alpha|^2}.
\]
所以 $\Delta'$ 在全体根集中不必对加法封闭。相应的加法封闭性体现在余根系中。$\alpha$ 的余根为 $\alpha^\vee=2\alpha/\langle\alpha,\alpha\rangle$，而 $\Delta'^\vee$ 在 $\Delta^\vee$ 中对加法封闭。引入某种对偶群 $G^\sim$，暂时还不是 $L$ 群。$G^\sim$ 为具有紧 Cartan 子群、根系为 $\Delta^\vee$ 的连通实群，安排成 $\Delta'^\vee$ 中的根为紧根，其他根为非紧根；要求其复化 $G_{\C}^\sim$ 无中心，Lie 代数为 $\mathfrak g^\sim$，Cartan 分解为 $\mathfrak g^\sim=\mathfrak k^\sim\oplus i\mathfrak p^\sim$。$W'$ 的元素给出 $\mathfrak k^\sim$ 的 Weyl 群元素，在 $G^\sim$ 中有代表元。$W$ 的元素给出 $\mathfrak g_{\C}^\sim$ 的 Weyl 群中保持紧根集稳定的元素，因而正规化 $\mathfrak k_{\C}^\sim$ 以及 $\mathfrak k^\sim\oplus i\mathfrak p^\sim$，进而正规化 $\mathfrak k^\sim,\mathfrak p^\sim$，从而正规化 $\mathfrak g^\sim$。由此 $R$ 嵌入商群 $\operatorname{Norm}_{G_{\C}^\sim}(\mathfrak g)/G^\sim$。这类论证可以进一步给出本情形下 $R=\sum\Z_2$ 的证明。

Langlands [14] 在这些想法的基础上，以 $L$ 群给出一般的 $W,W',R$ 表述。我们大致采用 Borel [2] 的记号，需要时遵照 Langlands [12] 的 $L$ 群构造。

先看上述例子与 $L$ 群的关系。当然，${}^LG^0$ 就是 $G_{\C}^\sim$（容许相差覆盖）。特征 $\xi$ 由各 $\xi(\gamma_\alpha)$ 指定，而 $\gamma_\alpha$ 的变换规律与余根相同，即 $\alpha^\vee$（所对应的特征）。知道 $\xi$ 在每个 $\alpha^\vee$ 上的值，就等于把 $\xi$ 看成对偶环面（dual torus）${}^LT^0$ 的元素。相应地，有 Weil 群到 ${}^LG$ 的同态 $\varphi$，由
\[
 \varphi(z)=1\ (z\in\C^\times),\qquad
 \varphi(\tau)=\xi\times\sigma\in{}^LT
\]
给出；这里 $\sigma$ 是 ${}^LG^0$ 的平凡外自同构（outer automorphism），因为 $G$ 分裂。中心化 $\varphi$ 的像 $\{1,\xi\times\sigma\}$ 的 Weyl 群元素就是 $W$ 的元素；这启发了下面的定理3.1。根向量被 $\varphi$ 的像中心化的根恰在 $\Delta'$ 中，这启发了定理3.3。上述 $R$ 与 $\operatorname{Norm}_{G_{\C}^\sim}(\mathfrak g)/G^\sim$ 同构的观察，也可用 $K^\sim$ 表述，从而启发定理3.4。

转到一般情形。在 $G$ 中固定极小抛物子群 $P_0=M_0A_0N_0$，令 $T$ 为包含 $A_0$ 的 Cartan 子群。设 ${}^LG^0\rtimes\{1,\sigma\}$ 是 $G$ 的 $L$ 群，带有极大环面 ${}^LT^0$ 及 Borel 子群 ${}^LB^0$。[2] 意义下的“标准相关抛物子群”（standard relevant parabolics）${}^LP$ 对应于 $G$ 的标准抛物子群 $P\supseteq P_0$。令 ${}^L(MA)$ 为 ${}^LP$ 的 Levi 分量（Levi component）。据 [12]，$G$ 的 $L$ 不可区分（$L$-indistinguishable）不可约容许表示的等价类，对应于某类 Weil 群到 ${}^LG$ 的同态 $\varphi$ 的（适当定义的）共轭类。表示缓增当且仅当 $\varphi$ 的像有界。它属于第2节中与 $P$ 对应的系列，当且仅当共轭以后 $\varphi$ 的像在 ${}^LP$ 的 Levi 分量 ${}^L(MA)$ 中，却不在其中任何更小的此类 Levi 分量中。以下可以假设 $\varphi(\C^\times)\subseteq{}^LT^0$，且 $\varphi(\tau)$ 在 ${}^LT^0$ 的正规化子中。本节末尾将给出例子，说明下面三个定理。

固定 $U_P(\xi,\nu)$ 中的一个直和项。相应的 $\varphi$ 对应 $G$ 中一个完整的表示 $L$ 类（$L$-class）。因 $\varphi$ 的像在 ${}^L(MA)$ 中，此 $L$ 类由从 $MA$ 的一个 $L$ 类诱导所得的那些表示组成。后者对应于将 $\varphi$ 视为取值在 ${}^L(MA)$ 中所得的映射 $\varphi^{MA}$，并由 $\nu$、$\xi$ 的中心特征及无穷小特征所刻画。而 $W,W',R$ 的定义只依赖于这三项数据。于是问题是如何直接用 $\varphi$ 定义 $W,W',R$。

\paragraph*{插话（Digression）。} 上述讨论及 D. Shelstad 的一个引理的简单推论是：$R$ 和对应于 $\varphi,\varphi^{MA}$ 的 $L$ 类的基数满足
\[
 |\varphi\text{ 的 }L\text{ 类}|=|R|\,|\varphi^{MA}\text{ 的 }L\text{ 类}|.
\]
事实上，$\varphi$ 的 $L$ 类中的表示，恰为所有 $U_P(\xi',\nu)$ 的不可约组成部分，其中 $\xi'\otimes\nu$ 属于 $\varphi^{MA}$ 的 $L$ 类。当 $\xi'$ 变化时，这些 $U_P(\xi',\nu)$ 的 $R$ 群相同。此外，Shelstad 引理指出：若 $\xi_1\otimes\nu$ 与 $\xi_2\otimes\nu$ 为 $L$ 等价但不等价，则 $U_P(\xi_1,\nu)$ 与 $U_P(\xi_2,\nu)$ 不相交。由此得到上述公式。

令 $S=\operatorname{Cent}(\operatorname{Image}\varphi)$，中心化子在 ${}^LG^0$ 中取；令 $S^0$ 为 $S$ 的单位连通分支，$\mathfrak s$ 为其 Lie 代数。

Weyl 群 $W(\mathfrak a)$ 嵌入 $({}^LG^0,{}^LT^0)$ 的 Weyl 群。具体说，$W(\mathfrak a)$ 的任一元素都有正规化 $\mathfrak m$ 的 Cartan 子代数并保持 $\mathfrak m$ 的正根的代表元；由此它给出 $({}^LG^0,{}^LT^0)$ 的 Weyl 群的唯一元素。在这一对应下，有：
\begin{originaltheorem}{定理 3.1（Langlands）}
$W=(\operatorname{Norm}({}^LT^0)\cap S)/({}^LT^0\cap S)$，其中正规化子在 ${}^LG^0$ 中取。
\end{originaltheorem}
将 $L$ 群数据中的格记为 $L,L^\vee$：
\[
 L=\Hom(T_{\C},\C^\times)\cong\Hom(\C^\times,{}^LT^0),
 \qquad
 L^\vee=\Hom({}^LT^0,\C^\times)\cong\Hom(\C^\times,T_{\C}).
\]
两个格之间的配对定义如下：若 $\lambda\in L=\Hom(\C^\times,{}^LT^0)$，$\lambda^\vee\in L^\vee=\Hom({}^LT^0,\C^\times)$，则 $\lambda^\vee\circ\lambda$ 是 $z\in\C^\times$ 的一个幂，$\langle\lambda,\lambda^\vee\rangle$ 表示此幂次。取值在 ${}^LT^0$ 的 $\varphi|_{\C^\times}$ 可形式地写成
\begin{equation}
 \varphi(z)=z^\mu\bar z^\nu=z^{\mu-\nu}(z\bar z)^\nu,
 \tag{3.1}
\end{equation}
其中 $\mu,\nu\in L\otimes\C$，且 $\mu-\nu\in L$。式 (3.1) 意味着，对所有 $\lambda^\vee\in L^\vee$，
\[
 \lambda^\vee(\varphi(z))=z^{\langle\mu-\nu,\lambda^\vee\rangle}
 (z\bar z)^{\langle\nu,\lambda^\vee\rangle}.
\]
$\varphi$ 的像不落在 $P$ 中任何真抛物子群中的条件意味着：恰对 ${}^L(MA)$ 的根，$\varphi(\tau)\alpha^\vee=-\alpha^\vee$，并且对 ${}^L(MA)$ 的所有根都有 $\langle\mu,\alpha^\vee\rangle\ne0$。（见 [12] 的引理3.3及引理3.1前一段。）

设 ${}^L\mathfrak t$ 为 ${}^LT$ 的 Lie 代数，将其按 $\varphi(\tau)$ 的 $+1,-1$ 特征空间分解为 ${}^L\mathfrak t={}^L\mathfrak t^+\oplus{}^L\mathfrak t^-$。
\begin{originaltheorem}{引理 3.2（Langlands）}
包含于 $\mathfrak s$ 中的 Lie 代数 ${}^L\mathfrak t^+$ 是 $\mathfrak s$ 的一个 Cartan 子代数。
\end{originaltheorem}
\begin{proof}
若 $X\in\mathfrak s$，条件 $X^{\varphi(\tau)}=X$ 意味着 $X$ 可写成 ${}^L\mathfrak t^+$ 中的一个元素与若干根向量之和 $Y$ 的和。此外 $Y$ 必须中心化 $\varphi(\C^\times)$。若 $Y=\sum Y_{\alpha^\vee}$，则
\[
 \sum Y_{\alpha^\vee}=Y=Y^{\varphi(z)}=
 \sum z^{\langle\mu-\nu,\alpha^\vee\rangle}(z\bar z)^{\langle\nu,\alpha^\vee\rangle}Y_{\alpha^\vee}.
\]
所以每个满足 $Y_{\alpha^\vee}\ne0$ 的 $\alpha^\vee$ 均有 $\langle\mu,\alpha^\vee\rangle=\langle\nu,\alpha^\vee\rangle=0$。因此
\[
 \mathfrak s\subseteq{}^L\mathfrak t^+
 +\sum_{\langle\mu,\alpha^\vee\rangle=\langle\nu,\alpha^\vee\rangle=0}\C X_{\alpha^\vee}.
\]
要证 ${}^L\mathfrak t^+$ 在 $\mathfrak s$ 中极大阿贝尔，只需说明这些 $X_{\alpha^\vee}$ 都不中心化 ${}^L\mathfrak t^+$，即相应 $\alpha^\vee$ 在 ${}^L\mathfrak t^+$ 上不为零。但如上所述，$\langle\mu,\alpha^\vee\rangle=0$ 蕴含 $\alpha^\vee$ 不是 ${}^L(MA)$ 的根，从而 $\varphi(\tau)\alpha^\vee\ne-\alpha^\vee$。因此 $\alpha^\vee$ 在 ${}^L\mathfrak t^+$ 上不为零。
\end{proof}
\begin{originaltheorem}{定理 3.3（Langlands）}
群 $W'$ 与 $(\mathfrak s,{}^L\mathfrak t^+)$ 的 Weyl 群典范同构。
\end{originaltheorem}
\begin{originaltheorem}{定理 3.4（Langlands）}
设 $G_{\C}$ 为半单单连通群。那么
\[
 F=({}^LT^0\cap S)/({}^LT^0\cap S^0)
\]
嵌入 $S/S^0$，商群与 $R$ 群同构。而且，若 $G$ 分裂且 $P$ 为极小抛物子群，则 $F=\{1\}$。
\end{originaltheorem}

\paragraph*{例。} 令 $G=\operatorname{SU}(2,1)$，为方便，以
\[
 g=\begin{pmatrix}2^{-1/2}&0&2^{-1/2}\\0&1&0\\-2^{-1/2}&0&2^{-1/2}\end{pmatrix}
\]
共轭之，从而可取\nobreak\footnote{\textbf{译者注（本次校订）：}本例所写环面的第三个对角元应为 $e^{-t+i\theta}$，从而 $T=\{\operatorname{diag}(e^{t+i\theta},e^{-2i\theta},e^{-t+i\theta})\}$。下文的 $H$ 也应为 $\operatorname{diag}(t+i\theta,-2i\theta,-t+i\theta)$，特征映射的输入矩阵须作相同更正。这样行列式才为一、$\operatorname{tr}H=0$，并且 $\langle\mu,H\rangle=t(c_1-c_2)+i\theta(c_1+c_2)$，与后文计算一致。}
\[
 T=\left\{\diag(e^{t+i\theta},e^{-2i\theta},e^{t-i\theta})\right\}.
\]
所用的 Borel 子群 $B$ 是上三角群。$G$ 拟分裂（quasi-split）但不分裂，$L$ 群为
\[
 {}^LG={}^LG^0\rtimes\Z_2=\operatorname{PGL}(3,\C)\rtimes\{1,\sigma\},
\]
其中 $\sigma$ 是 $\operatorname{PGL}(3,\C)$ 非平凡外自同构的一个特定实现。具体而言，${}^LT^0$ 为对角群，${}^LT={}^LT^0\rtimes\{1,\sigma\}$，${}^LB^0$ 为上三角群。于是得到标准简单余根，并在 $\operatorname{PGL}(3,\C)$ 中使用标准根向量；简单余根的根向量为
\[
 X_1=\begin{pmatrix}0&1&0\\0&0&0\\0&0&0\end{pmatrix},\qquad
 X_2=\begin{pmatrix}0&0&0\\0&0&1\\0&0&0\end{pmatrix},
\]
要求 $X_1^\sigma=X_2,X_2^\sigma=X_1$。显式地，$\sigma$ 在 ${}^LG^0$ 上的作用为
\[
 x^\sigma=w_0(x^{\mathrm{tr}})^{-1}w_0^{-1},\qquad
 w_0=\begin{pmatrix}0&0&1\\0&-1&0\\1&0&0\end{pmatrix},
\]
且 $\sigma^2=1$。在 Lie 代数 ${}^L\mathfrak t$ 上，有 $\diag(c_1,c_2,c_3)^\sigma=\diag(-c_3,-c_2,-c_1)$。

设 $\varphi$ 将 $\C^\times$ 映入 ${}^LT^0$，将 $\tau$ 映到 $w\times\sigma$，其中 $w\in\operatorname{Norm}({}^LT^0)$。容许共轭 $\varphi$ 后，$w$ 在 $\operatorname{Norm}({}^LT^0)/{}^LT^0$ 中只有两种可能，即 $1$ 与 $w_0$ 的陪集。无论哪种情形，$\varphi(\tau)$ 都给出 ${}^L\mathfrak t$ 的对合，具有 $+1,-1$ 特征空间 ${}^L\mathfrak t={}^L\mathfrak t^+\oplus{}^L\mathfrak t^-$。若 $[w]=[w_0]$，则 ${}^L\mathfrak t^+=0$，得到离散系列。我们关心相反情形 $[w]=[1]$。写 $w=x^\vee$，$x^\vee\in{}^LT^0$；此时 $\varphi(\tau)$ 在对角矩阵 ${}^LT^0$ 上的作用与 $\sigma$ 相同。

将 $L=\Hom(T_{\C},\C^\times)$ 的元素
\[
 \lambda(\diag(z,z^{-1}w^{-1},w))=z^aw^b
\]
简记为 $\lambda=a[z]+b[w]$。这里 $\{[z],[w]\}$ 为 $L$ 的一组基。令 $\{\delta_z,\delta_w\}$ 为 $L^\vee=\Hom({}^LT^0,\C^\times)$ 的对偶基。若
\begin{equation}
 \lambda^\vee(\diag(a_1,a_2,a_3)\bmod\text{中心})=(a_1/a_2)^c(a_3/a_2)^d,
 \tag{3.2}
\end{equation}
则 $\lambda^\vee=c\delta_z+d\delta_w$。$\varphi(\tau)$ 在 ${}^LT^0$ 上的作用导出 $L^\vee$ 的作用，再由对偶性导出 $L$ 上的作用，即
\[
 (a[z]+b[w])^{\varphi(\tau)}=-b[z]-a[w].
\]
对 $z\in\C^\times$，令 $\varphi(z)$ 如 (3.1)。因 Weil 群中有 $\bar z=\tau z\tau^{-1}$，得到 $\varphi(\bar z)=\varphi(z)^{\varphi(\tau)}$；由 (3.1) 推出 $\nu=\mu^{\varphi(\tau)}$。（这里 $\varphi(\tau)$ 在 $L\otimes\C$ 上仅作用于 $L$ 因子。）写
\[
 \mu=c_1[z]+c_2[w],\quad c_1,c_2\in\C;\qquad
 \nu=\mu^{\varphi(\tau)}=-c_2[z]-c_1[w].
\]
条件 $\mu-\nu\in L$ 意味着 $c_1+c_2\in\Z$。由 [12] 第27页可知，此时给主系列表示提供数据的 $T$ 特征为
\[
 e^H\longmapsto e^{\langle\mu,H+\sigma\bar H\rangle/2}=e^{\langle\mu,H\rangle}.
\]
若 $H=\diag(t+i\theta,-2i\theta,t-i\theta)$，得到特征
\[
 \diag(e^te^{i\theta},e^{-2i\theta},e^{t-i\theta})
 \longmapsto e^{t(c_1-c_2)}e^{i\theta(c_1+c_2)}.
\]
若 $\operatorname{Re}(c_1-c_2)>0$，得到非缓增表示。现在专门取 $c_1-c_2=0$ 的主系列，即取 $\mu=n[z]/2+n[w]/2$，它在 $M$ 部分对应 $e^{in\theta}$，在 $A$ 部分平凡。这些主系列可约（分裂成两部分）当且仅当 $n$ 为非零偶数。为求 $\varphi(z)$，将 ${}^LT^0$ 实现为 $\diag(a_1,a_2,a_3)\bmod\text{中心}$。在 (3.2) 中先取 $c=1,d=0$，再取 $c=0,d=1$，作用于 $\varphi(z)$，则
\[
 \begin{split}
 (a_1/a_2)(z)&=\delta_z\varphi(z)
 =z^{\langle\mu-\nu,\delta_z\rangle}(z\bar z)^{\langle\nu,\delta_z\rangle}\\
 &=z^{\langle2\mu,\delta_z\rangle}(z\bar z)^{\langle-\mu,\delta_z\rangle}=z^n/|z|^n,
 \end{split}
 \qquad
 (a_3/a_2)(z)=z^n/|z|^n.
\]
所以
\[
 \varphi(z)=\begin{pmatrix}
 z^n/|z|^n&0&0\\0&(z^n/|z|^n)^{-2}&0\\0&0&z^n/|z|^n
 \end{pmatrix}^{1/3}\bmod\text{中心}.
\]
这里相容地选取立方根；不同选择所造成的歧义在转到 $\operatorname{PGL}(3,\C)$ 时消失。

为完整计算 $\varphi$ 的像，还需确定 $\varphi(\tau)=x^\vee\times\sigma$。必须有 $\varphi(\tau)^2=\varphi(-1)$，故 $(x^\vee)(x^\vee)^\sigma=\varphi(-1)$。可取
\[
 x^\vee=\begin{cases}1,&n\text{ 为偶数},\\\diag(i,1,-i),&n\text{ 为奇数}.\end{cases}
\]
现在可计算本节定理涉及的 $S,S^0$。$\varphi$ 的像分三种不同情形：$n=0$、$n$ 奇数、$n$ 非零偶数。计算方式是先求 $S_1=\operatorname{Cent}\varphi(\C^\times)$，它作为环面的中心化子为连通群；然后在 $S_1$ 中求 $\varphi(\tau)$ 的中心化子。计算时尽量转到 Lie 代数上进行。对 $n=0$ 及 $n$ 奇数，使用如下技巧求 $S$：$S$ 包含在 ${}^LG^0$ 中 $\mathfrak s$ 的正规化子内，而计算得 $\mathfrak s=\mathfrak{sl}(2,\C)$。由于 $\mathfrak{sl}(2,\C)$ 没有外自同构，$S$ 对 $S^0$ 的共轭必为内自同构，故容易看出 $S=S^0$。结果如下：
\[
 \begin{array}{c|c|c}
 n&\mathfrak s&S\\\hline
 0&\left\{\begin{pmatrix}a&b&0\\c&0&b\\0&c&-a\end{pmatrix}\right\}&S^0\\[8pt]
 \text{奇数}&\left\{\begin{pmatrix}a&0&b\\0&0&0\\c&0&-a\end{pmatrix}\right\}&S^0\\[8pt]
 \text{非零偶数}&\left\{\begin{pmatrix}a&0&0\\0&0&0\\0&0&-a\end{pmatrix}\right\}&S^0\rtimes\{1,w_0\}
 \end{array}
\]
这些结果与本节定理及下列事实相符：
\[
 \begin{array}{c|cc}
 &W'&R\\\hline
 n=0\text{ 或奇数}&\Z_2&\{1\}\\
 n\text{ 非零偶数}&\{1\}&\Z_2
 \end{array}
\]

\origsection{4}{不可约缓增表示（Irreducible tempered representations）}
继续研究第2节的 $U_P(\xi,\nu)$，其中 $\xi$ 为 $M$ 的离散系列表示，$\nu$ 在 $\mathfrak a$ 上为纯虚。需要离散系列的参数化，我们采用 Harish-Chandra 的参数化。设 $T^-$ 为 $M$ 的一个紧 Cartan 子群，$\mathfrak t^-$ 为其 Lie 代数，$Z_M$ 为 $M$ 的中心。$M$ 的每个离散系列表示由 $i\mathfrak t^-$ 上的一个非奇异线性形式 $\lambda$ 及 $Z_M$ 上的一个特征 $\chi$ 决定。$\lambda$ 满足整性条件（integrality condition）：$\lambda-\rho$（$\rho$ 为某个次序下 $M$ 的正根半和）提升为 $T^-$ 的特征 $e^{\lambda-\rho}$；$\lambda,\chi$ 满足相容条件：$e^{\lambda-\rho}$ 与 $\chi$ 在公共定义域上一致。记 $\xi=\xi(\lambda,C,\chi)$，其中 $C$ 是 $i\mathfrak t^-$ 中使 $\lambda$ 占优的唯一 Weyl 室。两个此类 $\xi$ 等价当且仅当 $\chi=\chi'$，且在
\[
 W(T^-,M)=\operatorname{Norm}_M(T^-)/\operatorname{Cent}_M(T^-)
\]
中存在 $w$ 使 $w\lambda=\lambda'$、$wC=C'$。

[18] 讨论了“离散系列极限”（limits of discrete series）表示。它们可同样参数化，写作 $\xi=\xi(\lambda,C,\chi)$，只是允许 $\lambda$ 奇异，且使 $\lambda$ 占优的 $C$ 不止一种选择。这些表示或不可约，或为零；非零者的等价判据与离散系列相同。

基本表示（basic representation）是诱导表示 $U_P(\xi(\lambda,C,\chi),\nu)$，其中 $\xi(\lambda,C,\chi)$ 为 $M$ 上的离散系列极限。若 $\lambda$ 非奇异，故 $\xi$ 属于离散系列，便称此基本表示从离散系列诱导（induced from discrete series）。若对 $(\mathfrak m,\mathfrak t^-)$ 的每个根 $\alpha$，$\langle\lambda,\alpha\rangle=0$ 蕴含 $s_\alpha\notin W(T^-:M)$，则称基本表示具有非退化数据（nondegenerate data）。从离散系列诱导的表示自动具有非退化数据。

非退化条件有以下几个作用。首先，它排除了零表示。其次，它使基本表示的数据能够确定所用的抛物子群。例如，$\SL^\pm(2,\R)$ 的离散系列极限表示只有一个，它本身就是一个主系列表示。非退化条件要求将这个表示按主系列实现。一般地，它要求基本表示附属于尽可能小的抛物子群。第三，非退化性使第2节关于 $W,W',\Delta',R$ 的全部讨论推广到具有非退化数据的基本表示。[8]、[9] 的定理如下。
\begin{originaltheorem}{定理 4.1（Theorem 4.1）}
(a) 具有非退化数据的基本表示必为缓增，且不可约当且仅当其 $R$ 群平凡。

(b) 具有非退化数据的两个不可约基本表示 $U_P(\xi(\lambda,C,\chi),\nu)$ 与 $U_{P'}(\xi(\lambda',C',\chi'),\nu')$ 等价，当且仅当存在 $w\in G$，将 $M$ 送到 $M'$，$A$ 送到 $A'$，$\mathfrak t$ 送到 $\mathfrak t'$，且将 $(\lambda,C,\chi,\nu)$ 送到 $(\lambda',C',\chi',\nu')$。

(c) 每个不可约缓增表示都是基本表示，并可用非退化数据写出。
\end{originaltheorem}
$R$ 群可指出，将一个从离散系列诱导的表示分解时，会出现哪些具有非退化数据的基本表示。给定一个从离散系列诱导的表示，甚至从离散系列的整个一个 $L$ 类诱导的所有表示，它们的组成部分均为 $L$ 不可区分的。[8] 精确描述了如何从数据、$R$ 群及定理2.3(c)中的正交根集合读出此可约性。

\begin{thebibliography}{18}
\bibitem[1]{c05-1} J. Arthur, \textit{On the invariant distributions associated to weighted orbital integrals}, 预印本.
\bibitem[2]{c05-2} A. Borel, \textit{Automorphic L-functions}, Automorphic Forms, Representations, and L-functions, 预印本；本论文集第2册，27--61页.
\bibitem[3]{c05-3} F. Bruhat, \textit{Sur les représentations induites des groupes de Lie}, Bull. Soc. Math. France \textbf{84} (1956), 97--205.
\bibitem[4]{c05-4} Harish-Chandra, \textit{Harmonic analysis on real reductive groups. III}, Ann. of Math. (2) \textbf{104} (1976), 117--201.
\bibitem[5]{c05-5} A. W. Knapp, \textit{Commutativity of intertwining operators. II}, Bull. Amer. Math. Soc. \textbf{82} (1976), 271--273.
\bibitem[6]{c05-6} A. W. Knapp and E. M. Stein, \textit{Intertwining operators for semisimple groups}, Ann. of Math. (2) \textbf{93} (1971), 489--578.
\bibitem[7]{c05-7} A. W. Knapp and E. M. Stein, \textit{Singular integrals and the principal series. III, IV}, Proc. Nat. Acad. Sci. U.S.A. \textbf{71} (1974), 4622--4624; \textbf{72} (1975), 2459--2461.
\bibitem[8]{c05-8} A. W. Knapp and G. Zuckerman, \textit{Classification of irreducible tempered representations of semisimple Lie groups}, Proc. Nat. Acad. Sci. U.S.A. \textbf{73} (1976), 2178--2180.
\bibitem[9]{c05-9} A. W. Knapp and G. Zuckerman, \textit{Classification theorems for representations of semisimple Lie groups}, Non-Commutative Harmonic Analysis, Lecture Notes in Math., vol. 587, Springer-Verlag, New York, 1977, pp. 138--159.
\bibitem[10]{c05-10} R. A. Kunze and E. M. Stein, \textit{Uniformly bounded representations and harmonic analysis of the $2\times2$ real unimodular group}, Amer. J. Math. \textbf{82} (1960), 1--62.
\bibitem[11]{c05-11} R. A. Kunze and E. M. Stein, \textit{Uniformly bounded representations. III}, Amer. J. Math. \textbf{89} (1967), 385--442.
\bibitem[12]{c05-12} R. P. Langlands, \textit{On the classification of irreducible representations of real algebraic groups}, 油印讲义，Institute for Advanced Study, 1973.
\bibitem[13]{c05-13} R. P. Langlands, \textit{On the functional equations satisfied by Eisenstein series}, Lecture Notes in Math., vol. 544, Springer-Verlag, New York, 1976.
\bibitem[14]{c05-14} R. P. Langlands, \textit{Notes on the Knapp-Zuckerman theory}, 油印讲义，Institute for Advanced Study, Oct. 1976；1977年5月25日致 A. W. Knapp 的补充信.
\bibitem[15]{c05-15} G. Schiffmann, \textit{Intégrales d'entrelacement et fonctions de Whittaker}, Bull. Soc. Math. France \textbf{99} (1971), 3--72.
\bibitem[16]{c05-16} A. Silberger, \textit{The Knapp-Stein dimension theorem for p-adic groups}, 预印本.
\bibitem[17]{c05-17} N. R. Wallach, \textit{On Harish-Chandra's generalized C-functions}, Amer. J. Math. \textbf{97} (1975), 386--403.
\bibitem[18]{c05-18} G. Zuckerman, \textit{Tensor products of finite and infinite dimensional representations of semisimple Lie groups}, Ann. of Math. (2) \textbf{106} (1977), 295--308.
\end{thebibliography}
\noindent Cornell University（康奈尔大学）。\\
Yale University（耶鲁大学）。
